\section{Successive degrees and the remaining zero geometry}
\label{sec:zero-geometry}

The factorization already relates consecutive approximants exactly:
\begin{equation}
 \begin{aligned}
 \Delta_N(z)&:=G_{N+1}(z)-G_N(z)\\
 &=L_N(z)\left(E_N(z)+
       \frac{\beta_{N+1}z^2}{(2N+1)(2N+2)}E_{N+1}(z)\right).
 \end{aligned}
 \label{eq:degree-step}
\end{equation}
This follows by subtracting two copies of
\eqref{eq:error-factorization}. It identifies a signed correction;
the unknown error factors prevent treating it as a closed recurrence
for finite data alone.

On the common real sign range of Corollary~\ref{cor:brackets},
\(\Delta_N(x)\) has strict sign \((-1)^N\), for \(x\ne0\).
If \(r_1<r_2\) are two real zeros of \(G_N\) in that range,
\(G_{N+1}(r_1)\) and \(G_{N+1}(r_2)\) have the same nonzero sign.
Hence \(G_{N+1}\) cannot have exactly one simple zero and no other
zeros between them: a single simple zero would reverse that sign.
This explains why ordinary strict interlacing of consecutive Fourier
readouts is not supplied by the classical interlacing of their nodes.
It asserts neither that such endpoints exist nor a count of the
canonical readout zeros.

\subsection{What a phase comparison would need}

The simplest example is \(f(x)=\sin x\), \(g(x)=\cos x\).
The point \((g(x),f(x))\) travels around the unit circle, and
\(g(x)+\ii f(x)=\e^{\ii x}\) has continuously chosen angle \(x\).
Each passage through the horizontal axis is a zero of \(f\).
A positive angle derivative says the motion keeps turning in one
direction, so consecutive axis crossings give an ordered zero count.
Bounds on that derivative can also control the spacing of the crossings.
The angle of \(f\) alone jumps at its zeros; a companion supplies a
nonvanishing pair through which a continuous angle can be followed.

For real differentiable functions \(f,g\) that do not vanish
simultaneously on an interval, a continuous phase satisfies
\begin{equation}
 \frac{\dd}{\dd x}\arg\bigl(g(x)+\ii f(x)\bigr)
       =\frac{g(x)f'(x)-f(x)g'(x)}{f(x)^2+g(x)^2}.
 \label{eq:companion-phase}
\end{equation}
This follows by taking the imaginary part of
\((g'+\ii f')/(g+\ii f)\).
The numerator is a Wronskian, with the sign convention displayed.
For the polynomial pair, \eqref{eq:polynomial-wronskian} gives
its sign as a sum of squares. For \(f=G_N\), \(g=G_{N+1}\),
no corresponding positive identity has been established.
The preceding same-sign observation also prevents assuming an
ordinary strict interlacing law for this pair.

Thus a phase approach needs a justified companion and a proved
Wronskian sign or quantitative variation bound in the actual
observation variable. Such a statement would be additional
information about the readout, distinct from the small argument
of its error factor \(E_N\). Equation~\eqref{eq:companion-phase}
is background algebra, not a proposed theta inequality.

This derivative concerns movement in the observation variable at a
fixed degree. Following a particular zero as the degree increases is a
different question, and can be studied directly by local convergence.

\subsection{A local theorem explaining the role of slope}

The following elementary lemma makes precise the perturbative formula
used to motivate a possible source-specific estimate. A real entire
function here means \(f(\bar z)=\overline{f(z)}\).

\begin{lemma}[A simple-zero displacement bound]
\label{lem:zero-motion}
Let \(f,g\) be real entire functions and let \(r\in\R\) satisfy
\(f(r)=0\). Fix \(\rho>0\) and \(M_2,D\ge0\).
On the closed disk \(|z-r|\le\rho\), suppose
\begin{equation}
 |f'(r)|\ge s>0,\qquad |f''(z)|\le M_2,\qquad
 |g(z)-f(z)|\le D,\qquad
 D+\tfrac12M_2\rho^2<s\rho.
 \label{eq:local-zero-hypotheses}
\end{equation}
Then \(g\) has exactly one zero in this disk's interior, counted with
multiplicity. It is real and simple, and denoting it by \(r_g\),
\begin{equation}
 |r_g-r|\le\frac{D}{s-M_2\rho/2}.
 \label{eq:zero-displacement}
\end{equation}
For \(f_t=f+t(g-f)\), \(0\le t\le1\), its unique zero \(r(t)\)
in the disk is differentiable and
\begin{equation}
 r'(t)=-\frac{(g-f)(r(t))}{f_t'(r(t))},\qquad
 r'(0)=-\frac{g(r)-f(r)}{f'(r)}.
 \label{eq:zero-velocity}
\end{equation}
\end{lemma}

\begin{proof}
Taylor's integral remainder on a straight segment gives
\[
 |f(z)-f'(r)(z-r)|\le\tfrac12M_2|z-r|^2.
\]
On the boundary, the sum of this remainder and \(t(g-f)\) has
modulus strictly less than \(|f'(r)(z-r)|\). Rouch\'e's theorem
therefore gives exactly one zero of every \(f_t\) inside. Conjugation
symmetry makes it real, since a nonreal zero would bring a distinct
conjugate; the multiplicity count makes it simple.
At \(r_g\), Taylor's inequality also gives
\[
 D\ge|f(r_g)|\ge |r_g-r|\bigl(s-M_2|r_g-r|/2\bigr)
 \ge |r_g-r|(s-M_2\rho/2),
\]
proving \eqref{eq:zero-displacement}. For differentiability, the
unique simple zero can locally be written as the argument-principle
integral
\[
 r(t)=\frac{1}{2\pi\ii}\int_{|z-r|=\rho}
                  z\,\frac{f_t'(z)}{f_t(z)}\,\dd z.
\]
Nonvanishing on the contour persists in a complex neighborhood
of each real \(t\). Differentiation under the integral makes \(r(t)\)
holomorphic there. Differentiating \(f_t(r(t))=0\) proves
\eqref{eq:zero-velocity}.
\end{proof}

For \(f=G_N\), \(g=G_{N+1}\), the correction in this lemma is
\(\Delta_N\). A bound on it over a disk contained in \(|\Im z|\le Y\le1\)
follows from \eqref{eq:error-factorization} and
\eqref{eq:phase-bound}, for example
\[
 |\Delta_N(z)|\le
 |L_N(z)|(1+\epsilon_N(|z|,Y))
 +|L_{N+1}(z)|(1+\epsilon_{N+1}(|z|,Y)).
\]
Upper derivative bounds are also available from the finite sum:
if \(|\Im z|\le Y\) and the nodes have modulus at most \(\ell_N\),
then
\begin{equation}
 |G_N''(z)|\le\sum_j\omega_{j,N}x_{j,N}^2\e^{Y|x_{j,N}|}
                  \le \ell_N^2\e^{Y\ell_N}.
 \label{eq:finite-derivative-upper}
\end{equation}
The missing ingredient is a useful \emph{lower} slope or
cluster-separation estimate, together with enough information to
account for other zeros. Lemma~\ref{lem:zero-motion} is a conditional
local statement for any real entire functions. Its hypotheses have
not been proved for a growing canonical family of zero neighborhoods.
It imposes simplicity only within the particular disk; the overall
program allows multiple zeros and finite approximants with nonreal zeros.

\subsection{Following a fixed zero as degree increases}
\label{sec:fixed-zero}

Local uniform convergence already makes the idea of an older,
stabilized zero precise. The label is an isolating neighborhood of a
limiting zero, rather than the rank of an eigenvalue of \(J_N\).
The degree \(N\) is not the number of Fourier zeros: even \(G_2\)
has infinitely many. Increasing \(N\) does not add just one new
Fourier zero.

\begin{proposition}[Stability of an isolated zero cluster]
\label{prop:fixed-cluster}
Let \(r\in\C\) be a zero of \(F\) of multiplicity \(m\).
Choose \(\rho>0\) so that the closed disk \(D=\{|z-r|\le\rho\}\)
contains no other zero of \(F\). For all sufficiently large \(N\),
\(G_N\) has exactly \(m\) zeros in the interior of \(D\), counted
with multiplicity, and all of them tend to \(r\).
More explicitly, define the nonvanishing holomorphic factor
\(A(z)=F(z)/(z-r)^m\), with its removable value at \(r\), and put
\begin{equation}
 c=\min_{z\in D}|A(z)|>0,\qquad
 \eta_N=\sup_{z\in D}|F(z)-G_N(z)|.
 \quad
 G_N(w)=0,\ w\in D
 \ \Longrightarrow\
 |w-r|\le(\eta_N/c)^{1/m}.
 \label{eq:cluster-localization}
\end{equation}
If \(r\) is real and simple, the unique zero \(r_N\) in this disk
is real and simple for every sufficiently large \(N\).
\end{proposition}

\begin{proof}
The minimum of \(|F|\) on the boundary is positive.
Proposition~\ref{prop:compact} and Rouch\'e's theorem therefore give
the count \(m\) for all sufficiently large \(N\).
At a zero \(w\) of \(G_N\), factorization gives
\[
 c|w-r|^m\le|F(w)|=|F(w)-G_N(w)|\le\eta_N.
\]
Since \(\eta_N\to0\), this proves the localization and convergence.
For a real center the disk is invariant under conjugation.
When the total count is one, a nonreal root or a multiple root is
impossible.
\end{proof}

In a fixed bounded region with zero-free boundary, choose disjoint
small disks around its finitely many \(F\)-zeros. The compact
remainder has a positive minimum of \(|F|\), so eventually \(G_N\)
has no zeros there. Thus all zeros in that fixed region can be
matched by clusters. Individual labels are unique for simple zeros;
a multiple zero permits splitting and does not give a preferred
label for each constituent.

The complete error law gives more than qualitative stability at a
fixed simple real zero.

\begin{corollary}[Fixed-zero displacement and degree ratio]
\label{cor:fixed-zero-drift}
Let \(r\in\R\) be a fixed simple zero of \(F\), and let \(r_N\)
be the uniquely matched zeros above. Then \(r\ne0\), and
\begin{equation}
 \begin{aligned}
 r_N-r&=\frac{L_N(r)}{F'(r)}(1+o(1)),\\
 |r_N-r|&\le
 \frac{2|L_N(r)|}{|F'(r)|}
       \bigl(1+\epsilon_N(|r|,0)\bigr)
       \qquad\hbox{for all sufficiently large }N.
 \end{aligned}
 \label{eq:fixed-zero-displacement}
\end{equation}
Furthermore,
\begin{equation}
 \frac{r_{N+1}-r}{r_N-r}
 =-\frac{\beta_{N+1}r^2}{(2N+1)(2N+2)}(1+o(1))
 =O_r\!\left(\frac{\log^2(N+1)}{N^2}\right).
 \label{eq:fixed-zero-ratio}
\end{equation}
The denominator on the left is nonzero for all sufficiently large
\(N\). The displacements eventually alternate in sign, decrease
in absolute value at every step, and their even and odd subsequences
approach \(r\) monotonically from opposite sides.
\end{corollary}

\begin{proof}
We have \(F(0)=1\), so \(r\ne0\). Uniform convergence on a slightly
larger disk, and Cauchy's derivative formula, give \(G_N'\to F'\)
uniformly near \(r\). The exact error law evaluated at \(r\) gives
\(G_N(r)=-L_N(r)E_N(r)\). Consequently,
\[
 (r_N-r)\int_0^1G_N'\bigl(r+t(r_N-r)\bigr)\,\dd t
       =L_N(r)E_N(r).
\]
The integral tends to the nonzero number \(F'(r)\), while
Theorem~\ref{thm:phase} gives \(E_N(r)\to1\).
This proves the asymptotic and, once the integral's modulus is at
least \(|F'(r)|/2\), the displayed bound.
All factors are real, \(L_N(r)\) has sign \((-1)^N\), and
\(E_N(r)>0\) eventually. Taking consecutive ratios and using
\eqref{eq:leading-ratio} gives \eqref{eq:fixed-zero-ratio}.
The envelope bounds \(\beta_{N+1}=O(\log^2(N+1))\), so the
ratio is eventually negative with absolute value less than one.
The monotonicity assertions follow.
\end{proof}

The constants and starting degree depend on the chosen zero and
its isolation. This is a fixed-zero asymptotic, not a threshold
uniform over an increasing collection of zeros. To anchor a zero
without first knowing one of \(F\), one can instead isolate a single
zero of \(G_{N_0}\) in a conjugation-symmetric disk and prove
\(|F-G_{N_0}|<|G_{N_0}|\) on its boundary. Rouch\'e then certifies
one real simple zero of \(F\) there, to which the corollary applies.
Lemma~\ref{lem:zero-motion} gives one sufficient slope/curvature
test for this boundary comparison.

Stabilization alone does not establish reality: an isolated nonreal
zero of \(F\), if present, would also have a stable approximating
cluster by Proposition~\ref{prop:fixed-cluster}. A multiple real
zero can be approached by nonreal conjugate pairs. The established
statement concerns every fixed region; control on expanding regions
still requires the additional separation and count information.

\subsection{Clusters, resolution, and physical heat time}
\label{sec:cluster-resolution}

An unresolved pair need not be assigned two persistent individual labels.
It can instead be retained as a cluster with a certified total count.
The resolution here concerns what the approximation error allows one to
infer about \(F\), not an intrinsic inability to distinguish exact zeros
of a known \(G_N\). The scale depends on local separation and boundary
modulus as well as degree. In particular it is not a universal radius
depending only on \(N\).

There are two independent parameters. Increasing \(N\) improves the
approximation at a fixed physical heat time \(\lambda\); varying
\(\lambda\) changes the actual field. We now write \(G_{N,\lambda}\)
and \(F_\lambda\) explicitly. From \eqref{eq:normalized-field},
\begin{equation}
 \partial_\lambda F_\lambda
 =-\partial_z^2F_\lambda-m_2(\lambda)F_\lambda,
 \qquad m_2(\lambda)=\int u^2\,\dd\mu_\lambda(u).
 \label{eq:normalized-heat}
\end{equation}
The scalar normalization term does not move zeros. The unnormalized
\(H_\lambda\) satisfies \(\partial_\lambda H_\lambda=-\partial_z^2H_\lambda\).
Thus increasing \(\lambda\) is backward heat time; the reversed clock
\(\tau=1/2-\lambda\) has the opposite orientation.
The degree-dependent readouts are not assumed to obey that same exact
heat equation: their nodes and weights change with \(\lambda\).
The double-exponential tail permits differentiation on compact sets of
complex \((\lambda,z)\), making \(H\) jointly entire. Since
\(Z_\lambda>0\) on the real axis, the normalized field is jointly
holomorphic near each compact real time interval and spatial disk.

\begin{proposition}[Compact convergence uniform in heat time]
\label{prop:uniform-heat-compact}
For a compact real interval \(I=[\lambda_-,\lambda_+]\) and \(R<\infty\),
\begin{equation}
 \sup_{\lambda\in I,\,|z|\le R}|F_\lambda(z)-G_{N,\lambda}(z)|
 \le \frac{Z_{\lambda_+}}{Z_{\lambda_-}}
       \sum_{k=N}^\infty\frac{m_{2k}(\lambda_+)R^{2k}}{(2k)!}
 \longrightarrow0.
 \label{eq:uniform-heat-compact}
\end{equation}
\end{proposition}
\begin{proof}
Positivity gives \(Z_\lambda\ge Z_{\lambda_-}\) and
\(m_{2k}(\lambda)\le
(Z_{\lambda_+}/Z_{\lambda_-})m_{2k}(\lambda_+)\).
Apply \eqref{eq:compact-bound}; the dominating series converges by the
complete-source tail bound. All finite Jacobi coefficients, and hence
the readouts, depend continuously on real \(\lambda\), since their
moment determinants are positive and their moments are analytic.
\end{proof}

Fix a circle \(\Gamma=\{|z-a|=\rho\}\), \(a\in\R\), whose boundary
is zero-free for every \(\lambda\in I\). Suppose its interior contains
exactly two zeros of \(F_\lambda\), counted with multiplicity.
Such a circle and interval exist locally around any isolated double
zero. Set
\begin{equation}
 \eta=\min_{\lambda\in I,\,z\in\Gamma}|F_\lambda(z)|>0,
 \quad \varepsilon_N=\sup_{\lambda\in I,\,z\in\Gamma}
          |G_{N,\lambda}(z)-F_\lambda(z)|,
 \quad t_N=\varepsilon_N/\eta.
 \label{eq:cluster-boundary-budget}
\end{equation}
For \(t_N<1\), Rouch\'e gives exactly two zeros of \(G_{N,\lambda}\)
inside throughout \(I\), even when the pair cannot be separated.
The boundary minimum is a genuine hypothesis, not supplied by an upper
error bound alone. A finite-model alternative is given below.

Use contour moments to describe the cluster without choosing root labels:
\begin{equation}
 p_k(\lambda)=\frac{1}{2\pi\ii}\int_\Gamma
       (z-a)^k\frac{\partial_zF_\lambda(z)}{F_\lambda(z)}\,\dd z,
 \qquad
 c=a+\frac{p_1}{2},\qquad D=\frac{p_2}{2}-\frac{p_1^2}{4}.
 \label{eq:cluster-coordinates}
\end{equation}
Define \(p_{k,N},c_N,D_N\) by replacing \(F_\lambda\) with
\(G_{N,\lambda}\). The argument principle says \(p_0=2\), and the
residue theorem says \(p_k=(z_1-a)^k+(z_2-a)^k\). Therefore
\begin{equation}
 c=\frac{z_1+z_2}{2},\qquad
 D=\left(\frac{z_1-z_2}{2}\right)^2,\qquad
 \{z_1,z_2\}=\{c+\sqrt D,c-\sqrt D\}.
 \label{eq:pair-discriminant}
\end{equation}
The unordered pair is independent of the square-root branch. For real
\(\lambda\), conjugation symmetry makes \(c,D\) real. Positive \(D\)
means two distinct real zeros; negative \(D\) means a conjugate nonreal
pair; \(D=0\) means one real double zero. These symmetric coordinates
remain analytic in \(\lambda\) through the collision, as is immediate
from their contour definitions and joint analyticity of the field.

\begin{proposition}[Error bounds for an unresolved pair]
\label{prop:pair-resolution}
Under the preceding hypotheses and \(t_N<1\), put
\(b_N=-\log(1-t_N)\). Uniformly on \(I\),
\begin{equation}
 |p_{k,N}-p_k|\le k\rho^k b_N\quad(k\ge1),\qquad
 |c_N-c|\le\tfrac12\rho b_N,\qquad
 |D_N-D|\le 2\rho^2 b_N=:d_N.
 \label{eq:pair-resolution}
\end{equation}
In particular \(d_N\to0\). If \(D_N>d_N\), the limiting pair is
real and distinct; if \(D_N<-d_N\), it is nonreal. The band
\(|D_N|\le d_N\) is unresolved by this estimate and implies
\(|z_1-z_2|\le2\sqrt{2d_N}\).
\end{proposition}
\begin{proof}
On a neighborhood of \(\Gamma\), the ratio \(G_{N,\lambda}/F_\lambda\)
has the single-valued logarithm given by the convergent series for
\(\log(1+w)\), with \(w=(G_{N,\lambda}-F_\lambda)/F_\lambda\).
Its modulus on the circle is at most \(b_N\). Integration by parts
on the closed contour gives
\[
 p_{k,N}-p_k=-\frac{k}{2\pi\ii}\int_\Gamma
          (z-a)^{k-1}\log(G_{N,\lambda}/F_\lambda)\,\dd z.
\]
The circle length \(2\pi\rho\) proves the first bound. Both pairs
lie inside the circle, so \(|p_1|,|p_{1,N}|\le2\rho\).
Substitution into \eqref{eq:cluster-coordinates} gives the other two
bounds. The sign conclusions follow from reality of \(D,D_N\), and
the unresolved band gives \(|D|\le2d_N\).
\end{proof}

The decreasing moment-tail allowance in \eqref{eq:uniform-heat-compact}
may replace \(\varepsilon_N\), giving a decreasing resolution budget
once that allowance is below \(\eta\).

For a comparison using finite data, suppose instead that the readout
has exactly two zeros inside throughout \(I\), and
\(\gamma_N=\min_{I\times\Gamma}|G_{N,\lambda}|>\varepsilon_N\).
Rouch\'e transfers that count to \(F_\lambda\). Interchanging the
two functions in the logarithm proves the same bounds with
\(b_N=-\log(1-\varepsilon_N/\gamma_N)\). This version requires
the model boundary margin and count, but no prior knowledge of the
limiting pair's individual positions.

\begin{proposition}[Resolution near a fixed double collision]
\label{prop:double-collision-resolution}
Suppose \(F_{\lambda_*}\) has an isolated double zero at the real
point \(a\). For the local cluster coordinates just defined,
\begin{equation}
 D(\lambda_*)=0,\qquad D'(\lambda_*)=2,\qquad
 D(\lambda)=2(\lambda-\lambda_*)+O((\lambda-\lambda_*)^2).
 \label{eq:double-collision-slope}
\end{equation}
On a sufficiently small fixed interval \(I\) about \(\lambda_*\),
take \(D'\ge\kappa>0\). Every unresolved time satisfies
\begin{equation}
 |D_N(\lambda)|\le d_N
 \quad\Longrightarrow\quad
 |\lambda-\lambda_*|\le2d_N/\kappa.
 \label{eq:collision-time-resolution}
\end{equation}
For all sufficiently large \(N\), there is a time \(\lambda_N\in I\)
at which the finite readout has a double zero in the disk. Every such
time obeys \(|\lambda_N-\lambda_*|\le d_N/\kappa\).
Uniqueness of that finite-degree time is not asserted.
\end{proposition}
\begin{proof}
At \(\lambda_*\), write \(w=z-a\) and
\(F_{\lambda_*}(z)=w^2A(z)\), with \(A\) nonvanishing on the disk.
Then \(p_1=p_2=0\). Differentiate the contour formula and integrate
by parts. Equation~\eqref{eq:normalized-heat} gives
\[
 p_2'(\lambda_*)=-\frac{2}{2\pi\ii}\int_\Gamma
       w\frac{\partial_\lambda F}{F}\,\dd z
 =\frac{2}{2\pi\ii}\int_\Gamma w\frac{F_{zz}}{F}\,\dd z=4,
\]
because \(F_{zz}/F=2/w^2+4A'/(wA)+A''/A\). The normalization term
integrates to zero. Hence \(D'=p_2'/2=2\) at the event.
Analyticity proves the expansion and allows the choice of \(\kappa\).
The mean value theorem gives
\(|D(\lambda)|\ge\kappa|\lambda-\lambda_*|\), which proves
\eqref{eq:collision-time-resolution} and the stronger bound when
\(D_N=0\). Finally \(D\) has opposite signs at the two endpoints
of \(I\). Uniform convergence makes \(D_N\) have those signs too,
so continuity supplies a zero of \(D_N\).
\end{proof}

The elementary model \((z-a)^2-2(\lambda-\lambda_*)\) displays the
same orientation: as \(\lambda\) increases, a conjugate pair meets
and separates along the real axis. In reversed time the real pair
collides and leaves the axis. Its root separation is of square-root
size in elapsed time, while \(D\) is smooth. For a double zero the
spatial resolution is therefore of square-root size in the error,
whereas the time uncertainty is of linear size in the discriminant
error. This is a local approximation theorem about a fixed isolated
event. It neither asserts that a canonical collision exists nor
excludes one, and it is shared by colliding control fields with the
same local hypotheses. Higher multiplicities require more cluster
coordinates; no uniform conclusion on expanding height/time domains
is supplied here.

For background on individual real-zero dynamics in physical heat time,
Rodgers and Tao \cite[\S4, Theorem~11, Eq.~(56)]{rodgerstao} give a
principal-value interaction law in their stated range
\(\Lambda<\lambda\le0\). That simple-real-zero regime is distinct
from passing through a double event. The contour argument above needs
no individual label or division by the vanishing root separation.


\subsection{From fixed regions to the global objective}

The preceding cluster count requires no simplicity assumption.
To show that its multiplicity is accounted for by real zeros,
one still needs a real-zero count or confinement argument.
Tracking selected simple zeros does not rule out untracked complex pairs.

For completeness, the limiting objective can be stated precisely.
Suppose that the chosen \(G_N\) converge locally uniformly to the
nonzero entire \(F\), as Proposition~\ref{prop:compact} proves here.
Consider the assertion that there are
\(\mathcal R_N\to\infty\) and \(\delta_N\to0\) such that, for all
sufficiently large \(N\),
\begin{equation}
 G_N(z)=0,\quad |z|<\mathcal R_N
       \quad\Longrightarrow\quad |\Im z|\le\delta_N.
 \label{eq:qualitative-confinement}
\end{equation}
This assertion is equivalent to \(F\) having only real zeros.
One direction follows by choosing a small disk around a hypothetical
nonreal zero of \(F\), with zero-free boundary and disjoint from the
real axis. Uniform convergence and Rouch\'e then force zeros of
\(G_N\) in that disk, contradicting
\eqref{eq:qualitative-confinement}. Conversely, if \(F\) has only
real zeros, it has a positive minimum modulus on each compact set
\(\{|z|\le j,\ |\Im z|\ge1/j\}\). Uniform convergence makes \(G_N\)
zero-free there for all \(N\) above a threshold depending on \(j\).
A diagonal choice of increasing \(j=j(N)\) gives
\(\mathcal R_N=j(N)\), \(\delta_N=1/j(N)\).

For the unchanged time-zero theta field this is an RH-equivalent
reformulation, by its normalization in Section~\ref{sec:source}.
It is stated once as the objective, not as a new estimate. Prescribed
rates or an independently established degree invariant would be
additional mathematical content. A finite local estimate such as
Lemma~\ref{lem:zero-motion} is instead judged by its constants, domain,
and whether its hypotheses are actually available.
